\section{Introduction}
\label{sec:introduction}

Selecting the LoRA rank for each transformer layer is typically treated as a hyperparameter problem --- requiring grid search, gradient signals, or calibration runs before a single fine-tuning step begins. Yet the pretrained weight matrices that will be adapted already encode, in their singular value spectra, a structural fingerprint of how much representational complexity each layer developed during pretraining. If that fingerprint predicts how much rank a layer needs to absorb task-relevant signal, rank selection could be resolved in a single SVD pass --- before any training data is seen.

Low-Rank Adaptation (LoRA) \cite{Hu2022LoRA} achieves parameter-efficient fine-tuning by decomposing weight updates as $\Delta W = BA$ where $B \in \mathbb{R}^{d \times r}$, $A \in \mathbb{R}^{r \times k}$, and rank $r \ll \min(d,k)$. Its central weakness is the uniform rank assumption: every layer receives the same rank $r$, regardless of whether that layer's contribution to fine-tuning is large or small. This is known to be suboptimal: \citet{Aghajanyan2021Intrinsic} showed that intrinsic dimensionality varies substantially across transformer layers, and LAARA \cite{Tripathi2026LAARA} formally proves that uniform rank allocation is inferior to any Fisher-optimal per-layer assignment.

The deeper problem is that existing per-layer rank methods all require some form of training or calibration signal. Gradient-based methods --- AdaLoRA \cite{Zhang2023AdaLoRA}, DyLoRA \cite{Valipour2022DyLoRA}, La-LoRA \cite{Chen2025LaLoRA}, LAARA \cite{Tripathi2026LAARA} --- modify the rank during training using importance scores derived from weight updates. Calibration-based methods --- IFCLoRA \cite{Zhang2026IFCLoRA} --- run a forward pass over a calibration dataset before fine-tuning to estimate layer importance. Both families incur overhead: gradient-based methods add per-layer optimizer state and complexity; calibration methods require representative data and additional computation.

The gap is that no published work has tested whether a purely structural property of pretrained weights --- computed from $W_0$ alone with no training data --- can predict per-layer optimal LoRA rank with statistical significance. This gap exists because the field has focused on learned signals as the only reliable source of per-layer rank information, overlooking that pretraining itself shapes the singular value spectra in layer-type-specific ways.

\textbf{Our key insight:} The effective rank of a pretrained weight matrix --- $\text{erank}(W_0) = \exp(H(\sigma / \|\sigma\|_1))$ where $H$ is Shannon entropy and $\sigma$ are the singular values --- captures the same layer-complexity structure that an exhaustive oracle rank sweep independently discovers. Layers with high erank (spread-out spectra, many non-dominated directions) are precisely the layers the oracle assigns high LoRA rank to; layers with low erank (concentrated spectra, few dominant directions) receive low oracle rank. This connection is not incidental --- it reflects the geometry of how pretraining shapes each layer's representational capacity.

We introduce \textbf{erank-LoRA}, a zero-cost rank allocation strategy that computes $\text{erank}(W_0)$ for each weight matrix in a pretrained transformer and assigns LoRA rank proportionally, with no training data or calibration required. Our contributions are:

\begin{enumerate}
  \item \textbf{Empirical validation of erank-oracle correlation (P1):} We show that $\text{erank}(W_0)$ correlates with per-layer marginal PARA oracle ranks at Pearson $r = 0.984$ ($p = 0.0013$, $n=5$ oracle layers) for BERT-base-uncased --- far exceeding the pre-registered threshold of $r \geq 0.65$. FFN intermediate matrices (erank $\sim 704$--$710$) receive oracle rank $r=64$; attention layers including query and output projections (erank $\sim 557$--$590$) receive oracle rank $r=4$.

  \item \textbf{Participation ratio agreement (P5):} The participation ratio $\text{PR}(W_0)$ agrees with erank in layer ranking at Pearson $\rho = 0.968$, suggesting these two structural metrics capture the same underlying signal.

  \item \textbf{erank-oracle experimental protocol:} We define the PARA oracle sweep protocol and provide open implementations for erank computation, PARA oracle extraction, and correlation analysis, enabling replication and extension to additional model families.
\end{enumerate}

We organize the paper as follows: Section~\ref{sec:related} surveys related work on rank selection and structural weight analysis. Section~\ref{sec:methodology} describes our methodology. Section~\ref{sec:experiments} details our experimental setup. Section~\ref{sec:results} presents results for BERT-base-uncased. Section~\ref{sec:discussion} discusses implications, limitations, and future work. Section~\ref{sec:conclusion} concludes.
